\subsection{Dependencia estructural del funcional y normalización areal}

La construcción anterior sitúa el parámetro de Barbero--Immirzi en el
encuentro de tres determinaciones: incidencia marcada, orientación
dodecafásica y evaluación angular. La incidencia aporta los espacios
de sucesos; la orientación fija la cadena sobre esos espacios; la
evaluación transporta la cadena a una cantidad adimensional. Esta
descomposición permite conservar el origen de cada coeficiente al
utilizar después el funcional en la geometría de área.

La inclusión \(\mathcal F_6\hookrightarrow\mathcal F_8\) mantiene
el par marcado de la hexada. Su complemento tiene cardinal
\(2\binom62=30\). Este incremento reaparece como determinante del
mapa que reconstruye las ocupaciones de borde a partir del área y
del observable ponderado. La igualdad de ambos enteros posee aquí
una explicación concreta: los pesos sectoriales difieren precisamente
en el número de sucesos añadidos por la inclusión marcada.

\begin{proposition}[Incidencia y recuperabilidad de ocupaciones]
Sean \(n_6,n_8\) las ocupaciones agregadas de dos sectores y sean
\(r_6=|\mathcal F_6|\), \(r_8=|\mathcal F_8|\). El mapa
\[
 (n_6,n_8)\longmapsto(n_6+n_8,r_6n_6+r_8n_8)
\]
es invertible sobre su imagen cuando la inclusión añade al menos un
suceso. Para la inclusión construida, su determinante es 30.
\end{proposition}
\begin{proof}
La matriz del mapa es
\(\left(\begin{smallmatrix}1&1\\r_6&r_8\end{smallmatrix}\right)\).
Su determinante es \(r_8-r_6=|\mathcal F_8\setminus\mathcal F_6|\).
La inclusión considerada añade \(2\binom62=30\) sucesos. La inversión
lineal recupera ambas ocupaciones y conserva las restricciones enteras
propias de su imagen.
\end{proof}

El resultado recupera ocupaciones agregadas. La información sobre
cada enlace y el orden de los eventos permanece en el estado completo
y en sus lectores adicionales. Esta jerarquía da una formulación
precisa de la diferencia entre valor de área y genealogía de frontera:
el área conserva una suma; la pareja de observables conserva la
descomposición de esa suma por sectores; el registro conserva además
su distribución y su historia.

El parámetro \(\gamma_{\HMT}\) transporta esa incidencia a la escala
areal. Su interpretación estructural es la de un coeficiente orientado
del funcional de retorno hexádico--octádico. La realización gravitatoria
lo reconoce después en la acción de Holst. Así se distinguen su
definición intrínseca, su valor y su función en la ecuación física.
